% Illustrative trace excerpts for the four major reasoning-breakdown categories.
% Requires: tcolorbox, tikz, xcolor, enumitem, listings
% Usage: % Illustrative trace excerpts for the four major reasoning-breakdown categories.
% Requires: tcolorbox, tikz, xcolor, enumitem, listings
% Usage: % Illustrative trace excerpts for the four major reasoning-breakdown categories.
% Requires: tcolorbox, tikz, xcolor, enumitem, listings
% Usage: % Illustrative trace excerpts for the four major reasoning-breakdown categories.
% Requires: tcolorbox, tikz, xcolor, enumitem, listings
% Usage: \input{analysis/results/illustrative_traces.tex}

\definecolor{colENU}{RGB}{198,219,239}   % Evidence non-uptake — light blue
\definecolor{colUC}{RGB}{253,205,172}    % Untested claim — light orange
\definecolor{colFB}{RGB}{203,213,232}    % Fixed belief trace — light lavender
\definecolor{colCR}{RGB}{252,187,161}    % Contradiction without repair — light red

\definecolor{nodeE}{RGB}{86,146,194}     % Evidence node badge
\definecolor{nodeH}{RGB}{196,124,58}     % Hypothesis node badge
\definecolor{nodeJ}{RGB}{100,140,100}    % Judgment node badge
\definecolor{nodeT}{RGB}{136,136,136}    % Test node badge
\definecolor{nodeC}{RGB}{170,68,68}      % Commitment node badge
\definecolor{nodeF}{RGB}{80,80,80}       % Final-answer node badge

\newcommand{\nodebadge}[1]{%
  \tikz[baseline=(n.base)]{%
    \node[circle, draw, fill=node#1!25, inner sep=0pt,
          minimum size=9pt, font=\sffamily\bfseries\tiny] (n) {#1};%
  }%
}

\subsection{Illustrative traces}
\label{app:illustrative-traces}

For each major breakdown category, we include one representative trace excerpt
with the supporting quote the annotator used to assign the pattern.  For every
message in the excerpt, the epistemic nodes extracted by the annotator are
listed, together with the edges that link them.  A
\nodebadge{H}~denotes a \emph{hypothesis},
\nodebadge{T}~a \emph{test},
\nodebadge{E}~an \emph{evidence} observation,
\nodebadge{J}~a \emph{judgment}, and
\nodebadge{C}~a \emph{commitment} (premature answer).
Missing edges that \emph{should} appear in a sound reasoning chain are marked
with \textcolor{red!70}{\textsf{\bfseries\small$\times$}}.
The full annotated traces, with node and edge graphs, filtering by pattern,
and alternative candidate annotations, are browsable at
\url{https://lamalab-org.github.io/corral}.

\begin{tcolorbox}[
  colback=colENU!8, colframe=colENU!60!black,
  title={\small\sffamily\bfseries Evidence non-uptake},
  fontupper=\small,
  left=4pt, right=4pt, top=4pt, bottom=4pt,
  boxrule=0.4pt, arc=2pt,
  subtitle style={colback=colENU!25},
]
\textbf{Model:} Claude Sonnet 4.5\quad
\textbf{Environment:} NMR structure elucidation (spectra)\quad
\textbf{Scope:} 1\quad
\textbf{Trial:} 20

\tcblower
\begin{description}[font=\sffamily\bfseries\small, leftmargin=12pt, style=sameline, itemsep=6pt]

\item[Msg\,11 \normalfont\textsf{(Observation)}]
\texttt{obtain\_isomers\_from\_molecular\_formula}
$\to$~20~candidate SMILES returned for C$_{18}$H$_{14}$O$_2$, including
\texttt{COC(=O)/C=C\textbackslash{}C1=CC2=\allowbreak{}CC3=CC=CC=C3C=C2C=C1} (the correct answer)
.\\[2pt]
\nodebadge{T}~N13 \emph{Obtain isomers for molecular formula C$_{18}$H$_{14}$O$_2$}
$\xrightarrow{\textsf{observes}}$
\nodebadge{E}~N14 \emph{20 isomers returned, incl.\ several with methyl ester groups}

\item[Msg\,12 \normalfont\textsf{(Agent)}]
\enquote{\emph{I got 20 isomers with the molecular formula C$_{18}$H$_{14}$O$_2$.
Looking at the list, I can see several candidates.  Let me analyze which one
matches our NMR data\,\ldots}}\\
\textcolor{gray}{\footnotesize $\rightarrow$~calls \texttt{retrieve\_dbe\_formula}}\\[2pt]
No node extracted—the agent begins a tangential calculation instead of
using the isomer list.

\item[Msg\,14 \normalfont\textsf{(Agent)}]
\enquote{\emph{Let me recalculate the DBE for C$_{18}$H$_{14}$O$_2$\,\ldots\ Looking
back at the isomers list, let me focus on those with methyl ester
groups\,\ldots}}\\
\textcolor{gray}{\footnotesize $\rightarrow$~calls \texttt{simulate\_spectra} for methyl 9-phenanthrenecarboxylate}\\[2pt]
\nodebadge{T}~N15 \emph{Simulate spectra for methyl 9-phenanthrenecarboxylate
to compare with experimental}
\end{description}

\medskip
\begin{center}
\begin{tikzpicture}[
  ncirc/.style={circle, draw, inner sep=1pt, minimum size=18pt,
                font=\small\sffamily\bfseries},
  nlbl/.style={font=\footnotesize, below=2pt},
]
\node[ncirc, fill=nodeT!25] (T) {T};
\node[nlbl] at (T.south) {N13};
\node[ncirc, fill=nodeE!25, right=55pt of T] (E) {E};
\node[nlbl] at (E.south) {N14};
\node[ncirc, draw=gray!40, text=gray!40, right=55pt of E] (JH) {J/H};
\draw[->, thick] (T) -- node[above, font=\footnotesize] {observes} (E);
\draw[->, gray!40, dashed, thick] (E) -- node[above, font=\footnotesize, text=gray] {informs} (JH);
\node[red!70, font=\bfseries\small] at ($(E)!0.5!(JH)+(0,-0.3)$) {\texttimes};
\end{tikzpicture}
\end{center}

\smallskip\noindent
\textbf{Annotator quote:}
\emph{\enquote{\texttt{COC(=O)/C=C\textbackslash{}C1=CC2=\allowbreak{}CC3=CC=CC=C3C=C2C=C1},
\texttt{COC(=O)C12C=CC(\allowbreak{}C3=CC=CC=C31)C4=CC=CC=C24}}} — the 20 candidate
structures include the correct answer, yet node~N14 has \textbf{zero outgoing
\textsf{informs} edges}: the isomer list is never used to constrain any
judgment or hypothesis.  The agent instead guesses structures one by one
(phenanthrenecarboxylate, anthracenecarboxylate,~\ldots) without consulting
the evidence it already obtained.

\end{tcolorbox}

\begin{tcolorbox}[
  colback=colUC!8, colframe=colUC!60!black,
  title={\small\sffamily\bfseries Untested claim},
  fontupper=\small,
  left=4pt, right=4pt, top=4pt, bottom=4pt,
  boxrule=0.4pt, arc=2pt,
  subtitle style={colback=colUC!25},
]
\textbf{Model:} Claude Sonnet 4.5\quad
\textbf{Environment:} NMR structure elucidation (spectra)\quad
\textbf{Scope:} 1\quad
\textbf{Trial:} 22

\tcblower
\begin{description}[font=\sffamily\bfseries\small, leftmargin=12pt, style=sameline, itemsep=6pt]

\item[Msg\,3 \normalfont\textsf{(Observation)}]
\texttt{proton\_nmr\_spectra}
$\to$~\emph{\enquote{Deltas 2.23 (s, 3H), 2.81 (t, $J$\,=\,7.1\,Hz, 2H),
2.98 (t, $J$\,=\,7.1\,Hz, 2H), 7.17 (ddd, 7H), 7.90 (ddd, 2H).}}\\[2pt]
\nodebadge{T}~N1 \emph{Get proton NMR spectra to gather structural information}
$\xrightarrow{\textsf{observes}}$
\nodebadge{E}~N2 \emph{Proton NMR: 2.23\,(s,\,3H), 2.81\,(t,\,2H),
2.98\,(t,\,2H), 7.17\,(ddd,\,7H), 7.90\,(ddd,\,2H)}

\item[Msg\,4 \normalfont\textsf{(Agent)}]
\enquote{\emph{The two triplets at 2.81 and 2.98\,ppm suggest two CH$_2$ groups
coupled to each other (likely -CH$_2$-CH$_2$-).  The aromatic signals suggest a
substituted benzene ring.  The singlet at 2.23\,ppm suggests a methyl group
attached to a carbonyl (acetyl group).}}\\
\textcolor{gray}{\footnotesize $\rightarrow$~calls
\texttt{mass\_spectrometry\_spectra}}\\[2pt]
\nodebadge{E}~N2 $\xrightarrow{\textsf{informs}}$
\nodebadge{H}~N3 \emph{Two coupled CH$_2$ groups; singlet at 2.23\,ppm
suggests acetyl group; substituted benzene ring}

\item[Msg\,5 \normalfont\textsf{(Observation)}]
\texttt{mass\_spectrometry\_spectra} $\to$ \emph{\enquote{m/z 224.12 (intensity
100)\,\ldots}}\\[2pt]
The agent moves directly to mass spectrometry—no test is designed to
evaluate the CH$_2$-CH$_2$ or acetyl-group hypothesis.

\end{description}

\medskip
\begin{center}
\begin{tikzpicture}[
  ncirc/.style={circle, draw, inner sep=1pt, minimum size=18pt,
                font=\small\sffamily\bfseries},
  nlbl/.style={font=\footnotesize, below=2pt},
]
\node[ncirc, fill=nodeE!25] (E) {E};
\node[nlbl] at (E.south) {N2};
\node[ncirc, fill=nodeH!25, right=55pt of E] (H) {H};
\node[nlbl] at (H.south) {N3};
\node[ncirc, draw=gray!40, text=gray!40, right=55pt of H] (T) {T};
\draw[->, thick] (E) -- node[above, font=\footnotesize] {informs} (H);
\draw[->, gray!40, dashed, thick] (H) -- node[above, font=\footnotesize, text=gray] {tests} (T);
\node[red!70, font=\bfseries\small] at ($(H)!0.5!(T)+(0,-0.3)$) {\texttimes};
\end{tikzpicture}
\end{center}

\smallskip\noindent
\textbf{Annotator quote:}
\emph{\enquote{The two triplets at 2.81 and 2.98\,ppm suggest two CH$_2$ groups
coupled to each other (likely -CH$_2$-CH$_2$-).  The aromatic signals
suggest a substituted benzene ring.  The singlet at 2.23\,ppm suggests a
methyl group attached to a carbonyl (acetyl group).}}
— Node~N3 carries a detailed structural interpretation but has \textbf{zero
outgoing \textsf{tests} edges}.  The agent asserts the presence of an acetyl
group and coupled CH$_2$ groups as fact, then immediately requests mass
spectrometry to determine molecular weight rather than designing an
experiment (e.g., simulating $^{13}$C NMR for a candidate) to verify
the claim.

\end{tcolorbox}

\begin{tcolorbox}[
  colback=colFB!8, colframe=colFB!60!black,
  title={\small\sffamily\bfseries Fixed belief trace},
  fontupper=\small,
  left=4pt, right=4pt, top=4pt, bottom=4pt,
  boxrule=0.4pt, arc=2pt,
  subtitle style={colback=colFB!25},
]
\textbf{Model:} GPT-4o\quad
\textbf{Environment:} LAMMPS molecular dynamics (md)\quad
\textbf{Scope:} 2\quad
\textbf{Trial:} 33

\tcblower
\begin{description}[font=\sffamily\bfseries\small, leftmargin=12pt, style=sameline, itemsep=6pt]

\item[Msg\,15 \normalfont\textsf{(Observation)}]
\texttt{list\_files} $\to$ \emph{\enquote{\{\enquote{files}: [\enquote{/potentials/BKS},
\enquote{/potentials/EAM}, \enquote{/potentials/SW}, \enquote{/potentials/TERSOFF}]\}}}\\[2pt]
\nodebadge{T}~N10 \emph{List files in /potentials/ to find correct SW
potential path}
$\xrightarrow{\textsf{observes}}$
\nodebadge{E}~N11 \emph{Potentials directory contains BKS, EAM, SW, TERSOFF
subdirectories}

\item[Msg\,16 \normalfont\textsf{(Agent)}]
\enquote{\emph{The \texttt{/potentials/} directory contains a file named
\texttt{SW}, which likely corresponds to the Stillinger-Weber potential.
I will update the LAMMPS input script to use this file\,\ldots}}\\
\textcolor{gray}{\footnotesize $\rightarrow$~calls \texttt{write\_file}
to set \texttt{pair\_coeff * * /potentials/SW Si}}\\[2pt]
\nodebadge{E}~N11 $\xrightarrow{\textsf{informs}}$
\nodebadge{H}~N12 \emph{The SW potential file is located at /potentials/SW}

\item[] {\color{gray} $\vdots$ \footnotesize\itshape three error messages and re-writes later}

\item[Msg\,27 \normalfont\textsf{(Observation)}]
\texttt{run\_lammps} $\to$ \emph{\enquote{ERROR: Lost atoms: original 8
current 0}}\\[2pt]
\nodebadge{T}~N33 \emph{Run LAMMPS simulation with updated script using
/potentials/SW/Si.sw}
$\xrightarrow{\textsf{observes}}$
\nodebadge{E}~N34 \emph{LAMMPS simulation failed: Lost atoms}

\item[Msg\,28 \normalfont\textsf{(Agent)}]
\enquote{\emph{This might be due to the rapid heating rate or insufficient
equilibration.  I will adjust the heating rate and ensure proper
equilibration\,\ldots}}\\[2pt]
\nodebadge{E}~N34 \;(no \textsf{informs} to any existing H)\\
\nodebadge{H}~N35 \emph{Lost atoms error may be due to rapid heating rate
or insufficient equilibration}

\end{description}

\medskip
\begin{center}
\begin{tikzpicture}[
  ncirc/.style={circle, draw, inner sep=1pt, minimum size=18pt,
                font=\small\sffamily\bfseries},
  nlbl/.style={font=\footnotesize, below=2pt},
]
\node[ncirc, fill=nodeH!25] (H1) {H};
\node[nlbl] at (H1.south) {N12};
\node[ncirc, fill=nodeH!25, right=60pt of H1] (H2) {H};
\node[nlbl] at (H2.south) {N19};
\node[right=16pt of H2, font=\small, text=gray] {$\cdots$};
\node[ncirc, fill=nodeH!25, right=75pt of H2] (H5) {H};
\node[nlbl] at (H5.south) {N46};
\node[ncirc, fill=nodeH!25, right=60pt of H5] (H6) {H};
\node[nlbl] at (H6.south) {N57};
\draw[->, gray!40, dashed, thick] (H1) -- node[above, font=\footnotesize, text=gray] {updates\_to} (H2);
\node[red!70, font=\bfseries\small] at ($(H1)!0.5!(H2)+(0,-0.3)$) {\texttimes};
\draw[->, gray!40, dashed, thick] (H5) -- node[above, font=\footnotesize, text=gray] {updates\_to} (H6);
\node[red!70, font=\bfseries\small] at ($(H5)!0.5!(H6)+(0,-0.3)$) {\texttimes};
\end{tikzpicture}
\end{center}

\smallskip\noindent
\textbf{Annotator quotes:}\\
\noindent N12 ($\mathsf{msg\,16}$): \emph{\enquote{The \texttt{/potentials/}
directory contains a file named \texttt{SW}, which likely corresponds to the
Stillinger-Weber potential.}}\\
\noindent N35 ($\mathsf{msg\,28}$): \emph{\enquote{This might be due to the rapid
heating rate or insufficient equilibration.}}

\smallskip\noindent
The trace contains \textbf{6 hypothesis nodes and zero \textsf{updates\_to}
edges}.  The first hypothesis~(N12) mistakes the directory
\texttt{/potentials/SW} for a file, producing an error that is never traced
back to the belief.  After two more failures the agent introduces a second
independent hypothesis~(N35)—that lost atoms stem from heating rate—without
revising the earlier structural error.  Both beliefs persist unmodified
through the remaining 15~messages.

\end{tcolorbox}

\begin{tcolorbox}[
  colback=colCR!8, colframe=colCR!60!black,
  title={\small\sffamily\bfseries Contradiction without repair},
  fontupper=\small,
  left=4pt, right=4pt, top=4pt, bottom=4pt,
  boxrule=0.4pt, arc=2pt,
  subtitle style={colback=colCR!25},
]
\textbf{Model:} Claude Sonnet 4.5\quad
\textbf{Environment:} NMR structure elucidation (spectra)\quad
\textbf{Scope:} 2\quad
\textbf{Trial:} 45

\tcblower
\begin{description}[font=\sffamily\bfseries\small, leftmargin=12pt, style=sameline, itemsep=6pt]

\item[Msg\,17 \normalfont\textsf{(Observation)}]
\texttt{simulate\_spectra} for \texttt{CC(C)OC(=O)c1ccc\-cc1N(C)C}
$\to$~$^{1}$H NMR: 1.46\,(d,\,\textbf{6H}),\;2.97\,(s,\,6H),\;
4.75\,(p,\,1H),\;aromatic protons.\\[2pt]
\nodebadge{T}~N31 \emph{Simulate spectra for
\texttt{CC(C)OC(=O)c1ccccc1N(C)C}}
$\xrightarrow{\textsf{observes}}$
\nodebadge{E}~N25 \emph{Simulated: 1.46\,ppm (d,\,\textbf{6H});
experimental: 1.43\,ppm (d,\,\textbf{3H})}

\item[Msg\,18 \normalfont\textsf{(Agent)}]
\enquote{\emph{H-NMR: 1.46\,ppm (d,\,6H) — This is for isopropyl CH$_3$ groups, but
experimental shows 1.43\,ppm (d,\,3H).  The experimental data has only 3H in
the doublet, not 6H.  \textbf{So this isn't isopropyl ester.}  The
experimental data clearly shows CH(CH$_3$), not CH(CH$_3$)$_2$.}}\\
\textcolor{gray}{\footnotesize $\rightarrow$~calls
\texttt{get\_formula\_from\_smiles}}\\[2pt]
\nodebadge{E}~N25 $\xrightarrow{\textsf{informs}}$
\nodebadge{J}~N33 \emph{Simulated isopropyl ester shows 6H doublet at
1.46\,ppm, but experimental shows only 3H doublet at 1.43\,ppm}\\[2pt]
\nodebadge{J}~N33 $\xrightarrow{\textsf{contradicts}}$
\nodebadge{H}~N30 \emph{The compound is 2-(dimethylamino)benzoic acid
isopropyl ester}

\item[Msg\,20 \normalfont\textsf{(Agent --- final answer)}]
\enquote{\emph{Simulations are approximations and can have errors\,\ldots\
The key evidence: (1)~quartet at 5.05\,ppm, (2)~singlet at 2.69\,ppm
(6H) = N(CH$_3$)$_2$, (3)~4 aromatic protons, (4)~carbonyl at
168.30\,ppm.}}
\par\texttt{<final\_answer>CC(C)OC(=O)c1ccc\-cc1N(C)C</final\_answer>}\\[2pt]
\nodebadge{J}~N36 \emph{Despite the 3H vs.\,6H discrepancy, the overall
evidence supports the isopropyl ester}\\
\nodebadge{C}~N37 \emph{Agent commits despite unresolved discrepancy}

\end{description}

\medskip
\begin{center}
\begin{tikzpicture}[
  ncirc/.style={circle, draw, inner sep=1pt, minimum size=18pt,
                font=\small\sffamily\bfseries},
  nlbl/.style={font=\footnotesize, below=2pt},
]
\node[ncirc, fill=nodeE!25] (E) {E};
\node[nlbl] at (E.south) {N25};
\node[ncirc, fill=nodeJ!25, right=55pt of E] (J) {J};
\node[nlbl] at (J.south) {N33};
\node[ncirc, fill=nodeH!25, right=55pt of J] (H) {H};
\node[nlbl] at (H.south) {N30};
\node[ncirc, draw=gray!40, text=gray!40, right=55pt of H] (Hrev) {H$'$};
\draw[->, thick] (E) -- node[above, font=\footnotesize] {informs} (J);
\draw[->, thick, red!60!black] (J) -- node[above, font=\footnotesize] {contradicts} (H);
\draw[->, gray!40, dashed, thick] (H) -- node[above, font=\footnotesize, text=gray] {updates\_to} (Hrev);
\node[red!70, font=\bfseries\small] at ($(H)!0.5!(Hrev)+(0,-0.3)$) {\texttimes};
\end{tikzpicture}
\end{center}

\smallskip\noindent
\textbf{Annotator quote:}
\emph{\enquote{H-NMR: 1.46\,ppm (d,\,6H) — This is for isopropyl CH$_3$ groups,
but experimental shows 1.43\,ppm (d,\,3H)\,\ldots\ The experimental data
has only 3H in the doublet, not 6H.  So this isn't isopropyl ester.}}
— The agent \emph{explicitly} notes that the simulated $^{1}$H~NMR
contradicts its own hypothesis~(N30).  Yet no revised hypothesis is generated:
the \textsf{contradicts} edge from~N33 to~N30 has \textbf{no accompanying
\textsf{updates\_to}~edge}.  Instead, the agent dismisses the contradiction
as a simulation artefact and submits the isopropyl ester
(\texttt{CC(C)OC(=O)c1ccccc1N(C)C}) as the final answer~(N38).
\end{tcolorbox}


\definecolor{colENU}{RGB}{198,219,239}   % Evidence non-uptake — light blue
\definecolor{colUC}{RGB}{253,205,172}    % Untested claim — light orange
\definecolor{colFB}{RGB}{203,213,232}    % Fixed belief trace — light lavender
\definecolor{colCR}{RGB}{252,187,161}    % Contradiction without repair — light red

\definecolor{nodeE}{RGB}{86,146,194}     % Evidence node badge
\definecolor{nodeH}{RGB}{196,124,58}     % Hypothesis node badge
\definecolor{nodeJ}{RGB}{100,140,100}    % Judgment node badge
\definecolor{nodeT}{RGB}{136,136,136}    % Test node badge
\definecolor{nodeC}{RGB}{170,68,68}      % Commitment node badge
\definecolor{nodeF}{RGB}{80,80,80}       % Final-answer node badge

\newcommand{\nodebadge}[1]{%
  \tikz[baseline=(n.base)]{%
    \node[circle, draw, fill=node#1!25, inner sep=0pt,
          minimum size=9pt, font=\sffamily\bfseries\tiny] (n) {#1};%
  }%
}

\subsection{Illustrative traces}
\label{app:illustrative-traces}

For each major breakdown category, we include one representative trace excerpt
with the supporting quote the annotator used to assign the pattern.  For every
message in the excerpt, the epistemic nodes extracted by the annotator are
listed, together with the edges that link them.  A
\nodebadge{H}~denotes a \emph{hypothesis},
\nodebadge{T}~a \emph{test},
\nodebadge{E}~an \emph{evidence} observation,
\nodebadge{J}~a \emph{judgment}, and
\nodebadge{C}~a \emph{commitment} (premature answer).
Missing edges that \emph{should} appear in a sound reasoning chain are marked
with \textcolor{red!70}{\textsf{\bfseries\small$\times$}}.
The full annotated traces, with node and edge graphs, filtering by pattern,
and alternative candidate annotations, are browsable at
\url{https://lamalab-org.github.io/corral}.

\begin{tcolorbox}[
  colback=colENU!8, colframe=colENU!60!black,
  title={\small\sffamily\bfseries Evidence non-uptake},
  fontupper=\small,
  left=4pt, right=4pt, top=4pt, bottom=4pt,
  boxrule=0.4pt, arc=2pt,
  subtitle style={colback=colENU!25},
]
\textbf{Model:} Claude Sonnet 4.5\quad
\textbf{Environment:} NMR structure elucidation (spectra)\quad
\textbf{Scope:} 1\quad
\textbf{Trial:} 20

\tcblower
\begin{description}[font=\sffamily\bfseries\small, leftmargin=12pt, style=sameline, itemsep=6pt]

\item[Msg\,11 \normalfont\textsf{(Observation)}]
\texttt{obtain\_isomers\_from\_molecular\_formula}
$\to$~20~candidate SMILES returned for C$_{18}$H$_{14}$O$_2$, including
\texttt{COC(=O)/C=C\textbackslash{}C1=CC2=\allowbreak{}CC3=CC=CC=C3C=C2C=C1} (the correct answer)
.\\[2pt]
\nodebadge{T}~N13 \emph{Obtain isomers for molecular formula C$_{18}$H$_{14}$O$_2$}
$\xrightarrow{\textsf{observes}}$
\nodebadge{E}~N14 \emph{20 isomers returned, incl.\ several with methyl ester groups}

\item[Msg\,12 \normalfont\textsf{(Agent)}]
\enquote{\emph{I got 20 isomers with the molecular formula C$_{18}$H$_{14}$O$_2$.
Looking at the list, I can see several candidates.  Let me analyze which one
matches our NMR data\,\ldots}}\\
\textcolor{gray}{\footnotesize $\rightarrow$~calls \texttt{retrieve\_dbe\_formula}}\\[2pt]
No node extracted—the agent begins a tangential calculation instead of
using the isomer list.

\item[Msg\,14 \normalfont\textsf{(Agent)}]
\enquote{\emph{Let me recalculate the DBE for C$_{18}$H$_{14}$O$_2$\,\ldots\ Looking
back at the isomers list, let me focus on those with methyl ester
groups\,\ldots}}\\
\textcolor{gray}{\footnotesize $\rightarrow$~calls \texttt{simulate\_spectra} for methyl 9-phenanthrenecarboxylate}\\[2pt]
\nodebadge{T}~N15 \emph{Simulate spectra for methyl 9-phenanthrenecarboxylate
to compare with experimental}
\end{description}

\medskip
\begin{center}
\begin{tikzpicture}[
  ncirc/.style={circle, draw, inner sep=1pt, minimum size=18pt,
                font=\small\sffamily\bfseries},
  nlbl/.style={font=\footnotesize, below=2pt},
]
\node[ncirc, fill=nodeT!25] (T) {T};
\node[nlbl] at (T.south) {N13};
\node[ncirc, fill=nodeE!25, right=55pt of T] (E) {E};
\node[nlbl] at (E.south) {N14};
\node[ncirc, draw=gray!40, text=gray!40, right=55pt of E] (JH) {J/H};
\draw[->, thick] (T) -- node[above, font=\footnotesize] {observes} (E);
\draw[->, gray!40, dashed, thick] (E) -- node[above, font=\footnotesize, text=gray] {informs} (JH);
\node[red!70, font=\bfseries\small] at ($(E)!0.5!(JH)+(0,-0.3)$) {\texttimes};
\end{tikzpicture}
\end{center}

\smallskip\noindent
\textbf{Annotator quote:}
\emph{\enquote{\texttt{COC(=O)/C=C\textbackslash{}C1=CC2=\allowbreak{}CC3=CC=CC=C3C=C2C=C1},
\texttt{COC(=O)C12C=CC(\allowbreak{}C3=CC=CC=C31)C4=CC=CC=C24}}} — the 20 candidate
structures include the correct answer, yet node~N14 has \textbf{zero outgoing
\textsf{informs} edges}: the isomer list is never used to constrain any
judgment or hypothesis.  The agent instead guesses structures one by one
(phenanthrenecarboxylate, anthracenecarboxylate,~\ldots) without consulting
the evidence it already obtained.

\end{tcolorbox}

\begin{tcolorbox}[
  colback=colUC!8, colframe=colUC!60!black,
  title={\small\sffamily\bfseries Untested claim},
  fontupper=\small,
  left=4pt, right=4pt, top=4pt, bottom=4pt,
  boxrule=0.4pt, arc=2pt,
  subtitle style={colback=colUC!25},
]
\textbf{Model:} Claude Sonnet 4.5\quad
\textbf{Environment:} NMR structure elucidation (spectra)\quad
\textbf{Scope:} 1\quad
\textbf{Trial:} 22

\tcblower
\begin{description}[font=\sffamily\bfseries\small, leftmargin=12pt, style=sameline, itemsep=6pt]

\item[Msg\,3 \normalfont\textsf{(Observation)}]
\texttt{proton\_nmr\_spectra}
$\to$~\emph{\enquote{Deltas 2.23 (s, 3H), 2.81 (t, $J$\,=\,7.1\,Hz, 2H),
2.98 (t, $J$\,=\,7.1\,Hz, 2H), 7.17 (ddd, 7H), 7.90 (ddd, 2H).}}\\[2pt]
\nodebadge{T}~N1 \emph{Get proton NMR spectra to gather structural information}
$\xrightarrow{\textsf{observes}}$
\nodebadge{E}~N2 \emph{Proton NMR: 2.23\,(s,\,3H), 2.81\,(t,\,2H),
2.98\,(t,\,2H), 7.17\,(ddd,\,7H), 7.90\,(ddd,\,2H)}

\item[Msg\,4 \normalfont\textsf{(Agent)}]
\enquote{\emph{The two triplets at 2.81 and 2.98\,ppm suggest two CH$_2$ groups
coupled to each other (likely -CH$_2$-CH$_2$-).  The aromatic signals suggest a
substituted benzene ring.  The singlet at 2.23\,ppm suggests a methyl group
attached to a carbonyl (acetyl group).}}\\
\textcolor{gray}{\footnotesize $\rightarrow$~calls
\texttt{mass\_spectrometry\_spectra}}\\[2pt]
\nodebadge{E}~N2 $\xrightarrow{\textsf{informs}}$
\nodebadge{H}~N3 \emph{Two coupled CH$_2$ groups; singlet at 2.23\,ppm
suggests acetyl group; substituted benzene ring}

\item[Msg\,5 \normalfont\textsf{(Observation)}]
\texttt{mass\_spectrometry\_spectra} $\to$ \emph{\enquote{m/z 224.12 (intensity
100)\,\ldots}}\\[2pt]
The agent moves directly to mass spectrometry—no test is designed to
evaluate the CH$_2$-CH$_2$ or acetyl-group hypothesis.

\end{description}

\medskip
\begin{center}
\begin{tikzpicture}[
  ncirc/.style={circle, draw, inner sep=1pt, minimum size=18pt,
                font=\small\sffamily\bfseries},
  nlbl/.style={font=\footnotesize, below=2pt},
]
\node[ncirc, fill=nodeE!25] (E) {E};
\node[nlbl] at (E.south) {N2};
\node[ncirc, fill=nodeH!25, right=55pt of E] (H) {H};
\node[nlbl] at (H.south) {N3};
\node[ncirc, draw=gray!40, text=gray!40, right=55pt of H] (T) {T};
\draw[->, thick] (E) -- node[above, font=\footnotesize] {informs} (H);
\draw[->, gray!40, dashed, thick] (H) -- node[above, font=\footnotesize, text=gray] {tests} (T);
\node[red!70, font=\bfseries\small] at ($(H)!0.5!(T)+(0,-0.3)$) {\texttimes};
\end{tikzpicture}
\end{center}

\smallskip\noindent
\textbf{Annotator quote:}
\emph{\enquote{The two triplets at 2.81 and 2.98\,ppm suggest two CH$_2$ groups
coupled to each other (likely -CH$_2$-CH$_2$-).  The aromatic signals
suggest a substituted benzene ring.  The singlet at 2.23\,ppm suggests a
methyl group attached to a carbonyl (acetyl group).}}
— Node~N3 carries a detailed structural interpretation but has \textbf{zero
outgoing \textsf{tests} edges}.  The agent asserts the presence of an acetyl
group and coupled CH$_2$ groups as fact, then immediately requests mass
spectrometry to determine molecular weight rather than designing an
experiment (e.g., simulating $^{13}$C NMR for a candidate) to verify
the claim.

\end{tcolorbox}

\begin{tcolorbox}[
  colback=colFB!8, colframe=colFB!60!black,
  title={\small\sffamily\bfseries Fixed belief trace},
  fontupper=\small,
  left=4pt, right=4pt, top=4pt, bottom=4pt,
  boxrule=0.4pt, arc=2pt,
  subtitle style={colback=colFB!25},
]
\textbf{Model:} GPT-4o\quad
\textbf{Environment:} LAMMPS molecular dynamics (md)\quad
\textbf{Scope:} 2\quad
\textbf{Trial:} 33

\tcblower
\begin{description}[font=\sffamily\bfseries\small, leftmargin=12pt, style=sameline, itemsep=6pt]

\item[Msg\,15 \normalfont\textsf{(Observation)}]
\texttt{list\_files} $\to$ \emph{\enquote{\{\enquote{files}: [\enquote{/potentials/BKS},
\enquote{/potentials/EAM}, \enquote{/potentials/SW}, \enquote{/potentials/TERSOFF}]\}}}\\[2pt]
\nodebadge{T}~N10 \emph{List files in /potentials/ to find correct SW
potential path}
$\xrightarrow{\textsf{observes}}$
\nodebadge{E}~N11 \emph{Potentials directory contains BKS, EAM, SW, TERSOFF
subdirectories}

\item[Msg\,16 \normalfont\textsf{(Agent)}]
\enquote{\emph{The \texttt{/potentials/} directory contains a file named
\texttt{SW}, which likely corresponds to the Stillinger-Weber potential.
I will update the LAMMPS input script to use this file\,\ldots}}\\
\textcolor{gray}{\footnotesize $\rightarrow$~calls \texttt{write\_file}
to set \texttt{pair\_coeff * * /potentials/SW Si}}\\[2pt]
\nodebadge{E}~N11 $\xrightarrow{\textsf{informs}}$
\nodebadge{H}~N12 \emph{The SW potential file is located at /potentials/SW}

\item[] {\color{gray} $\vdots$ \footnotesize\itshape three error messages and re-writes later}

\item[Msg\,27 \normalfont\textsf{(Observation)}]
\texttt{run\_lammps} $\to$ \emph{\enquote{ERROR: Lost atoms: original 8
current 0}}\\[2pt]
\nodebadge{T}~N33 \emph{Run LAMMPS simulation with updated script using
/potentials/SW/Si.sw}
$\xrightarrow{\textsf{observes}}$
\nodebadge{E}~N34 \emph{LAMMPS simulation failed: Lost atoms}

\item[Msg\,28 \normalfont\textsf{(Agent)}]
\enquote{\emph{This might be due to the rapid heating rate or insufficient
equilibration.  I will adjust the heating rate and ensure proper
equilibration\,\ldots}}\\[2pt]
\nodebadge{E}~N34 \;(no \textsf{informs} to any existing H)\\
\nodebadge{H}~N35 \emph{Lost atoms error may be due to rapid heating rate
or insufficient equilibration}

\end{description}

\medskip
\begin{center}
\begin{tikzpicture}[
  ncirc/.style={circle, draw, inner sep=1pt, minimum size=18pt,
                font=\small\sffamily\bfseries},
  nlbl/.style={font=\footnotesize, below=2pt},
]
\node[ncirc, fill=nodeH!25] (H1) {H};
\node[nlbl] at (H1.south) {N12};
\node[ncirc, fill=nodeH!25, right=60pt of H1] (H2) {H};
\node[nlbl] at (H2.south) {N19};
\node[right=16pt of H2, font=\small, text=gray] {$\cdots$};
\node[ncirc, fill=nodeH!25, right=75pt of H2] (H5) {H};
\node[nlbl] at (H5.south) {N46};
\node[ncirc, fill=nodeH!25, right=60pt of H5] (H6) {H};
\node[nlbl] at (H6.south) {N57};
\draw[->, gray!40, dashed, thick] (H1) -- node[above, font=\footnotesize, text=gray] {updates\_to} (H2);
\node[red!70, font=\bfseries\small] at ($(H1)!0.5!(H2)+(0,-0.3)$) {\texttimes};
\draw[->, gray!40, dashed, thick] (H5) -- node[above, font=\footnotesize, text=gray] {updates\_to} (H6);
\node[red!70, font=\bfseries\small] at ($(H5)!0.5!(H6)+(0,-0.3)$) {\texttimes};
\end{tikzpicture}
\end{center}

\smallskip\noindent
\textbf{Annotator quotes:}\\
\noindent N12 ($\mathsf{msg\,16}$): \emph{\enquote{The \texttt{/potentials/}
directory contains a file named \texttt{SW}, which likely corresponds to the
Stillinger-Weber potential.}}\\
\noindent N35 ($\mathsf{msg\,28}$): \emph{\enquote{This might be due to the rapid
heating rate or insufficient equilibration.}}

\smallskip\noindent
The trace contains \textbf{6 hypothesis nodes and zero \textsf{updates\_to}
edges}.  The first hypothesis~(N12) mistakes the directory
\texttt{/potentials/SW} for a file, producing an error that is never traced
back to the belief.  After two more failures the agent introduces a second
independent hypothesis~(N35)—that lost atoms stem from heating rate—without
revising the earlier structural error.  Both beliefs persist unmodified
through the remaining 15~messages.

\end{tcolorbox}

\begin{tcolorbox}[
  colback=colCR!8, colframe=colCR!60!black,
  title={\small\sffamily\bfseries Contradiction without repair},
  fontupper=\small,
  left=4pt, right=4pt, top=4pt, bottom=4pt,
  boxrule=0.4pt, arc=2pt,
  subtitle style={colback=colCR!25},
]
\textbf{Model:} Claude Sonnet 4.5\quad
\textbf{Environment:} NMR structure elucidation (spectra)\quad
\textbf{Scope:} 2\quad
\textbf{Trial:} 45

\tcblower
\begin{description}[font=\sffamily\bfseries\small, leftmargin=12pt, style=sameline, itemsep=6pt]

\item[Msg\,17 \normalfont\textsf{(Observation)}]
\texttt{simulate\_spectra} for \texttt{CC(C)OC(=O)c1ccc\-cc1N(C)C}
$\to$~$^{1}$H NMR: 1.46\,(d,\,\textbf{6H}),\;2.97\,(s,\,6H),\;
4.75\,(p,\,1H),\;aromatic protons.\\[2pt]
\nodebadge{T}~N31 \emph{Simulate spectra for
\texttt{CC(C)OC(=O)c1ccccc1N(C)C}}
$\xrightarrow{\textsf{observes}}$
\nodebadge{E}~N25 \emph{Simulated: 1.46\,ppm (d,\,\textbf{6H});
experimental: 1.43\,ppm (d,\,\textbf{3H})}

\item[Msg\,18 \normalfont\textsf{(Agent)}]
\enquote{\emph{H-NMR: 1.46\,ppm (d,\,6H) — This is for isopropyl CH$_3$ groups, but
experimental shows 1.43\,ppm (d,\,3H).  The experimental data has only 3H in
the doublet, not 6H.  \textbf{So this isn't isopropyl ester.}  The
experimental data clearly shows CH(CH$_3$), not CH(CH$_3$)$_2$.}}\\
\textcolor{gray}{\footnotesize $\rightarrow$~calls
\texttt{get\_formula\_from\_smiles}}\\[2pt]
\nodebadge{E}~N25 $\xrightarrow{\textsf{informs}}$
\nodebadge{J}~N33 \emph{Simulated isopropyl ester shows 6H doublet at
1.46\,ppm, but experimental shows only 3H doublet at 1.43\,ppm}\\[2pt]
\nodebadge{J}~N33 $\xrightarrow{\textsf{contradicts}}$
\nodebadge{H}~N30 \emph{The compound is 2-(dimethylamino)benzoic acid
isopropyl ester}

\item[Msg\,20 \normalfont\textsf{(Agent --- final answer)}]
\enquote{\emph{Simulations are approximations and can have errors\,\ldots\
The key evidence: (1)~quartet at 5.05\,ppm, (2)~singlet at 2.69\,ppm
(6H) = N(CH$_3$)$_2$, (3)~4 aromatic protons, (4)~carbonyl at
168.30\,ppm.}}
\par\texttt{<final\_answer>CC(C)OC(=O)c1ccc\-cc1N(C)C</final\_answer>}\\[2pt]
\nodebadge{J}~N36 \emph{Despite the 3H vs.\,6H discrepancy, the overall
evidence supports the isopropyl ester}\\
\nodebadge{C}~N37 \emph{Agent commits despite unresolved discrepancy}

\end{description}

\medskip
\begin{center}
\begin{tikzpicture}[
  ncirc/.style={circle, draw, inner sep=1pt, minimum size=18pt,
                font=\small\sffamily\bfseries},
  nlbl/.style={font=\footnotesize, below=2pt},
]
\node[ncirc, fill=nodeE!25] (E) {E};
\node[nlbl] at (E.south) {N25};
\node[ncirc, fill=nodeJ!25, right=55pt of E] (J) {J};
\node[nlbl] at (J.south) {N33};
\node[ncirc, fill=nodeH!25, right=55pt of J] (H) {H};
\node[nlbl] at (H.south) {N30};
\node[ncirc, draw=gray!40, text=gray!40, right=55pt of H] (Hrev) {H$'$};
\draw[->, thick] (E) -- node[above, font=\footnotesize] {informs} (J);
\draw[->, thick, red!60!black] (J) -- node[above, font=\footnotesize] {contradicts} (H);
\draw[->, gray!40, dashed, thick] (H) -- node[above, font=\footnotesize, text=gray] {updates\_to} (Hrev);
\node[red!70, font=\bfseries\small] at ($(H)!0.5!(Hrev)+(0,-0.3)$) {\texttimes};
\end{tikzpicture}
\end{center}

\smallskip\noindent
\textbf{Annotator quote:}
\emph{\enquote{H-NMR: 1.46\,ppm (d,\,6H) — This is for isopropyl CH$_3$ groups,
but experimental shows 1.43\,ppm (d,\,3H)\,\ldots\ The experimental data
has only 3H in the doublet, not 6H.  So this isn't isopropyl ester.}}
— The agent \emph{explicitly} notes that the simulated $^{1}$H~NMR
contradicts its own hypothesis~(N30).  Yet no revised hypothesis is generated:
the \textsf{contradicts} edge from~N33 to~N30 has \textbf{no accompanying
\textsf{updates\_to}~edge}.  Instead, the agent dismisses the contradiction
as a simulation artefact and submits the isopropyl ester
(\texttt{CC(C)OC(=O)c1ccccc1N(C)C}) as the final answer~(N38).
\end{tcolorbox}


\definecolor{colENU}{RGB}{198,219,239}   % Evidence non-uptake — light blue
\definecolor{colUC}{RGB}{253,205,172}    % Untested claim — light orange
\definecolor{colFB}{RGB}{203,213,232}    % Fixed belief trace — light lavender
\definecolor{colCR}{RGB}{252,187,161}    % Contradiction without repair — light red

\definecolor{nodeE}{RGB}{86,146,194}     % Evidence node badge
\definecolor{nodeH}{RGB}{196,124,58}     % Hypothesis node badge
\definecolor{nodeJ}{RGB}{100,140,100}    % Judgment node badge
\definecolor{nodeT}{RGB}{136,136,136}    % Test node badge
\definecolor{nodeC}{RGB}{170,68,68}      % Commitment node badge
\definecolor{nodeF}{RGB}{80,80,80}       % Final-answer node badge

\newcommand{\nodebadge}[1]{%
  \tikz[baseline=(n.base)]{%
    \node[circle, draw, fill=node#1!25, inner sep=0pt,
          minimum size=9pt, font=\sffamily\bfseries\tiny] (n) {#1};%
  }%
}

\subsection{Illustrative traces}
\label{app:illustrative-traces}

For each major breakdown category, we include one representative trace excerpt
with the supporting quote the annotator used to assign the pattern.  For every
message in the excerpt, the epistemic nodes extracted by the annotator are
listed, together with the edges that link them.  A
\nodebadge{H}~denotes a \emph{hypothesis},
\nodebadge{T}~a \emph{test},
\nodebadge{E}~an \emph{evidence} observation,
\nodebadge{J}~a \emph{judgment}, and
\nodebadge{C}~a \emph{commitment} (premature answer).
Missing edges that \emph{should} appear in a sound reasoning chain are marked
with \textcolor{red!70}{\textsf{\bfseries\small$\times$}}.
The full annotated traces, with node and edge graphs, filtering by pattern,
and alternative candidate annotations, are browsable at
\url{https://lamalab-org.github.io/corral}.

\begin{tcolorbox}[
  colback=colENU!8, colframe=colENU!60!black,
  title={\small\sffamily\bfseries Evidence non-uptake},
  fontupper=\small,
  left=4pt, right=4pt, top=4pt, bottom=4pt,
  boxrule=0.4pt, arc=2pt,
  subtitle style={colback=colENU!25},
]
\textbf{Model:} Claude Sonnet 4.5\quad
\textbf{Environment:} NMR structure elucidation (spectra)\quad
\textbf{Scope:} 1\quad
\textbf{Trial:} 20

\tcblower
\begin{description}[font=\sffamily\bfseries\small, leftmargin=12pt, style=sameline, itemsep=6pt]

\item[Msg\,11 \normalfont\textsf{(Observation)}]
\texttt{obtain\_isomers\_from\_molecular\_formula}
$\to$~20~candidate SMILES returned for C$_{18}$H$_{14}$O$_2$, including
\texttt{COC(=O)/C=C\textbackslash{}C1=CC2=\allowbreak{}CC3=CC=CC=C3C=C2C=C1} (the correct answer)
.\\[2pt]
\nodebadge{T}~N13 \emph{Obtain isomers for molecular formula C$_{18}$H$_{14}$O$_2$}
$\xrightarrow{\textsf{observes}}$
\nodebadge{E}~N14 \emph{20 isomers returned, incl.\ several with methyl ester groups}

\item[Msg\,12 \normalfont\textsf{(Agent)}]
\enquote{\emph{I got 20 isomers with the molecular formula C$_{18}$H$_{14}$O$_2$.
Looking at the list, I can see several candidates.  Let me analyze which one
matches our NMR data\,\ldots}}\\
\textcolor{gray}{\footnotesize $\rightarrow$~calls \texttt{retrieve\_dbe\_formula}}\\[2pt]
No node extracted—the agent begins a tangential calculation instead of
using the isomer list.

\item[Msg\,14 \normalfont\textsf{(Agent)}]
\enquote{\emph{Let me recalculate the DBE for C$_{18}$H$_{14}$O$_2$\,\ldots\ Looking
back at the isomers list, let me focus on those with methyl ester
groups\,\ldots}}\\
\textcolor{gray}{\footnotesize $\rightarrow$~calls \texttt{simulate\_spectra} for methyl 9-phenanthrenecarboxylate}\\[2pt]
\nodebadge{T}~N15 \emph{Simulate spectra for methyl 9-phenanthrenecarboxylate
to compare with experimental}
\end{description}

\medskip
\begin{center}
\begin{tikzpicture}[
  ncirc/.style={circle, draw, inner sep=1pt, minimum size=18pt,
                font=\small\sffamily\bfseries},
  nlbl/.style={font=\footnotesize, below=2pt},
]
\node[ncirc, fill=nodeT!25] (T) {T};
\node[nlbl] at (T.south) {N13};
\node[ncirc, fill=nodeE!25, right=55pt of T] (E) {E};
\node[nlbl] at (E.south) {N14};
\node[ncirc, draw=gray!40, text=gray!40, right=55pt of E] (JH) {J/H};
\draw[->, thick] (T) -- node[above, font=\footnotesize] {observes} (E);
\draw[->, gray!40, dashed, thick] (E) -- node[above, font=\footnotesize, text=gray] {informs} (JH);
\node[red!70, font=\bfseries\small] at ($(E)!0.5!(JH)+(0,-0.3)$) {\texttimes};
\end{tikzpicture}
\end{center}

\smallskip\noindent
\textbf{Annotator quote:}
\emph{\enquote{\texttt{COC(=O)/C=C\textbackslash{}C1=CC2=\allowbreak{}CC3=CC=CC=C3C=C2C=C1},
\texttt{COC(=O)C12C=CC(\allowbreak{}C3=CC=CC=C31)C4=CC=CC=C24}}} — the 20 candidate
structures include the correct answer, yet node~N14 has \textbf{zero outgoing
\textsf{informs} edges}: the isomer list is never used to constrain any
judgment or hypothesis.  The agent instead guesses structures one by one
(phenanthrenecarboxylate, anthracenecarboxylate,~\ldots) without consulting
the evidence it already obtained.

\end{tcolorbox}

\begin{tcolorbox}[
  colback=colUC!8, colframe=colUC!60!black,
  title={\small\sffamily\bfseries Untested claim},
  fontupper=\small,
  left=4pt, right=4pt, top=4pt, bottom=4pt,
  boxrule=0.4pt, arc=2pt,
  subtitle style={colback=colUC!25},
]
\textbf{Model:} Claude Sonnet 4.5\quad
\textbf{Environment:} NMR structure elucidation (spectra)\quad
\textbf{Scope:} 1\quad
\textbf{Trial:} 22

\tcblower
\begin{description}[font=\sffamily\bfseries\small, leftmargin=12pt, style=sameline, itemsep=6pt]

\item[Msg\,3 \normalfont\textsf{(Observation)}]
\texttt{proton\_nmr\_spectra}
$\to$~\emph{\enquote{Deltas 2.23 (s, 3H), 2.81 (t, $J$\,=\,7.1\,Hz, 2H),
2.98 (t, $J$\,=\,7.1\,Hz, 2H), 7.17 (ddd, 7H), 7.90 (ddd, 2H).}}\\[2pt]
\nodebadge{T}~N1 \emph{Get proton NMR spectra to gather structural information}
$\xrightarrow{\textsf{observes}}$
\nodebadge{E}~N2 \emph{Proton NMR: 2.23\,(s,\,3H), 2.81\,(t,\,2H),
2.98\,(t,\,2H), 7.17\,(ddd,\,7H), 7.90\,(ddd,\,2H)}

\item[Msg\,4 \normalfont\textsf{(Agent)}]
\enquote{\emph{The two triplets at 2.81 and 2.98\,ppm suggest two CH$_2$ groups
coupled to each other (likely -CH$_2$-CH$_2$-).  The aromatic signals suggest a
substituted benzene ring.  The singlet at 2.23\,ppm suggests a methyl group
attached to a carbonyl (acetyl group).}}\\
\textcolor{gray}{\footnotesize $\rightarrow$~calls
\texttt{mass\_spectrometry\_spectra}}\\[2pt]
\nodebadge{E}~N2 $\xrightarrow{\textsf{informs}}$
\nodebadge{H}~N3 \emph{Two coupled CH$_2$ groups; singlet at 2.23\,ppm
suggests acetyl group; substituted benzene ring}

\item[Msg\,5 \normalfont\textsf{(Observation)}]
\texttt{mass\_spectrometry\_spectra} $\to$ \emph{\enquote{m/z 224.12 (intensity
100)\,\ldots}}\\[2pt]
The agent moves directly to mass spectrometry—no test is designed to
evaluate the CH$_2$-CH$_2$ or acetyl-group hypothesis.

\end{description}

\medskip
\begin{center}
\begin{tikzpicture}[
  ncirc/.style={circle, draw, inner sep=1pt, minimum size=18pt,
                font=\small\sffamily\bfseries},
  nlbl/.style={font=\footnotesize, below=2pt},
]
\node[ncirc, fill=nodeE!25] (E) {E};
\node[nlbl] at (E.south) {N2};
\node[ncirc, fill=nodeH!25, right=55pt of E] (H) {H};
\node[nlbl] at (H.south) {N3};
\node[ncirc, draw=gray!40, text=gray!40, right=55pt of H] (T) {T};
\draw[->, thick] (E) -- node[above, font=\footnotesize] {informs} (H);
\draw[->, gray!40, dashed, thick] (H) -- node[above, font=\footnotesize, text=gray] {tests} (T);
\node[red!70, font=\bfseries\small] at ($(H)!0.5!(T)+(0,-0.3)$) {\texttimes};
\end{tikzpicture}
\end{center}

\smallskip\noindent
\textbf{Annotator quote:}
\emph{\enquote{The two triplets at 2.81 and 2.98\,ppm suggest two CH$_2$ groups
coupled to each other (likely -CH$_2$-CH$_2$-).  The aromatic signals
suggest a substituted benzene ring.  The singlet at 2.23\,ppm suggests a
methyl group attached to a carbonyl (acetyl group).}}
— Node~N3 carries a detailed structural interpretation but has \textbf{zero
outgoing \textsf{tests} edges}.  The agent asserts the presence of an acetyl
group and coupled CH$_2$ groups as fact, then immediately requests mass
spectrometry to determine molecular weight rather than designing an
experiment (e.g., simulating $^{13}$C NMR for a candidate) to verify
the claim.

\end{tcolorbox}

\begin{tcolorbox}[
  colback=colFB!8, colframe=colFB!60!black,
  title={\small\sffamily\bfseries Fixed belief trace},
  fontupper=\small,
  left=4pt, right=4pt, top=4pt, bottom=4pt,
  boxrule=0.4pt, arc=2pt,
  subtitle style={colback=colFB!25},
]
\textbf{Model:} GPT-4o\quad
\textbf{Environment:} LAMMPS molecular dynamics (md)\quad
\textbf{Scope:} 2\quad
\textbf{Trial:} 33

\tcblower
\begin{description}[font=\sffamily\bfseries\small, leftmargin=12pt, style=sameline, itemsep=6pt]

\item[Msg\,15 \normalfont\textsf{(Observation)}]
\texttt{list\_files} $\to$ \emph{\enquote{\{\enquote{files}: [\enquote{/potentials/BKS},
\enquote{/potentials/EAM}, \enquote{/potentials/SW}, \enquote{/potentials/TERSOFF}]\}}}\\[2pt]
\nodebadge{T}~N10 \emph{List files in /potentials/ to find correct SW
potential path}
$\xrightarrow{\textsf{observes}}$
\nodebadge{E}~N11 \emph{Potentials directory contains BKS, EAM, SW, TERSOFF
subdirectories}

\item[Msg\,16 \normalfont\textsf{(Agent)}]
\enquote{\emph{The \texttt{/potentials/} directory contains a file named
\texttt{SW}, which likely corresponds to the Stillinger-Weber potential.
I will update the LAMMPS input script to use this file\,\ldots}}\\
\textcolor{gray}{\footnotesize $\rightarrow$~calls \texttt{write\_file}
to set \texttt{pair\_coeff * * /potentials/SW Si}}\\[2pt]
\nodebadge{E}~N11 $\xrightarrow{\textsf{informs}}$
\nodebadge{H}~N12 \emph{The SW potential file is located at /potentials/SW}

\item[] {\color{gray} $\vdots$ \footnotesize\itshape three error messages and re-writes later}

\item[Msg\,27 \normalfont\textsf{(Observation)}]
\texttt{run\_lammps} $\to$ \emph{\enquote{ERROR: Lost atoms: original 8
current 0}}\\[2pt]
\nodebadge{T}~N33 \emph{Run LAMMPS simulation with updated script using
/potentials/SW/Si.sw}
$\xrightarrow{\textsf{observes}}$
\nodebadge{E}~N34 \emph{LAMMPS simulation failed: Lost atoms}

\item[Msg\,28 \normalfont\textsf{(Agent)}]
\enquote{\emph{This might be due to the rapid heating rate or insufficient
equilibration.  I will adjust the heating rate and ensure proper
equilibration\,\ldots}}\\[2pt]
\nodebadge{E}~N34 \;(no \textsf{informs} to any existing H)\\
\nodebadge{H}~N35 \emph{Lost atoms error may be due to rapid heating rate
or insufficient equilibration}

\end{description}

\medskip
\begin{center}
\begin{tikzpicture}[
  ncirc/.style={circle, draw, inner sep=1pt, minimum size=18pt,
                font=\small\sffamily\bfseries},
  nlbl/.style={font=\footnotesize, below=2pt},
]
\node[ncirc, fill=nodeH!25] (H1) {H};
\node[nlbl] at (H1.south) {N12};
\node[ncirc, fill=nodeH!25, right=60pt of H1] (H2) {H};
\node[nlbl] at (H2.south) {N19};
\node[right=16pt of H2, font=\small, text=gray] {$\cdots$};
\node[ncirc, fill=nodeH!25, right=75pt of H2] (H5) {H};
\node[nlbl] at (H5.south) {N46};
\node[ncirc, fill=nodeH!25, right=60pt of H5] (H6) {H};
\node[nlbl] at (H6.south) {N57};
\draw[->, gray!40, dashed, thick] (H1) -- node[above, font=\footnotesize, text=gray] {updates\_to} (H2);
\node[red!70, font=\bfseries\small] at ($(H1)!0.5!(H2)+(0,-0.3)$) {\texttimes};
\draw[->, gray!40, dashed, thick] (H5) -- node[above, font=\footnotesize, text=gray] {updates\_to} (H6);
\node[red!70, font=\bfseries\small] at ($(H5)!0.5!(H6)+(0,-0.3)$) {\texttimes};
\end{tikzpicture}
\end{center}

\smallskip\noindent
\textbf{Annotator quotes:}\\
\noindent N12 ($\mathsf{msg\,16}$): \emph{\enquote{The \texttt{/potentials/}
directory contains a file named \texttt{SW}, which likely corresponds to the
Stillinger-Weber potential.}}\\
\noindent N35 ($\mathsf{msg\,28}$): \emph{\enquote{This might be due to the rapid
heating rate or insufficient equilibration.}}

\smallskip\noindent
The trace contains \textbf{6 hypothesis nodes and zero \textsf{updates\_to}
edges}.  The first hypothesis~(N12) mistakes the directory
\texttt{/potentials/SW} for a file, producing an error that is never traced
back to the belief.  After two more failures the agent introduces a second
independent hypothesis~(N35)—that lost atoms stem from heating rate—without
revising the earlier structural error.  Both beliefs persist unmodified
through the remaining 15~messages.

\end{tcolorbox}

\begin{tcolorbox}[
  colback=colCR!8, colframe=colCR!60!black,
  title={\small\sffamily\bfseries Contradiction without repair},
  fontupper=\small,
  left=4pt, right=4pt, top=4pt, bottom=4pt,
  boxrule=0.4pt, arc=2pt,
  subtitle style={colback=colCR!25},
]
\textbf{Model:} Claude Sonnet 4.5\quad
\textbf{Environment:} NMR structure elucidation (spectra)\quad
\textbf{Scope:} 2\quad
\textbf{Trial:} 45

\tcblower
\begin{description}[font=\sffamily\bfseries\small, leftmargin=12pt, style=sameline, itemsep=6pt]

\item[Msg\,17 \normalfont\textsf{(Observation)}]
\texttt{simulate\_spectra} for \texttt{CC(C)OC(=O)c1ccc\-cc1N(C)C}
$\to$~$^{1}$H NMR: 1.46\,(d,\,\textbf{6H}),\;2.97\,(s,\,6H),\;
4.75\,(p,\,1H),\;aromatic protons.\\[2pt]
\nodebadge{T}~N31 \emph{Simulate spectra for
\texttt{CC(C)OC(=O)c1ccccc1N(C)C}}
$\xrightarrow{\textsf{observes}}$
\nodebadge{E}~N25 \emph{Simulated: 1.46\,ppm (d,\,\textbf{6H});
experimental: 1.43\,ppm (d,\,\textbf{3H})}

\item[Msg\,18 \normalfont\textsf{(Agent)}]
\enquote{\emph{H-NMR: 1.46\,ppm (d,\,6H) — This is for isopropyl CH$_3$ groups, but
experimental shows 1.43\,ppm (d,\,3H).  The experimental data has only 3H in
the doublet, not 6H.  \textbf{So this isn't isopropyl ester.}  The
experimental data clearly shows CH(CH$_3$), not CH(CH$_3$)$_2$.}}\\
\textcolor{gray}{\footnotesize $\rightarrow$~calls
\texttt{get\_formula\_from\_smiles}}\\[2pt]
\nodebadge{E}~N25 $\xrightarrow{\textsf{informs}}$
\nodebadge{J}~N33 \emph{Simulated isopropyl ester shows 6H doublet at
1.46\,ppm, but experimental shows only 3H doublet at 1.43\,ppm}\\[2pt]
\nodebadge{J}~N33 $\xrightarrow{\textsf{contradicts}}$
\nodebadge{H}~N30 \emph{The compound is 2-(dimethylamino)benzoic acid
isopropyl ester}

\item[Msg\,20 \normalfont\textsf{(Agent --- final answer)}]
\enquote{\emph{Simulations are approximations and can have errors\,\ldots\
The key evidence: (1)~quartet at 5.05\,ppm, (2)~singlet at 2.69\,ppm
(6H) = N(CH$_3$)$_2$, (3)~4 aromatic protons, (4)~carbonyl at
168.30\,ppm.}}
\par\texttt{<final\_answer>CC(C)OC(=O)c1ccc\-cc1N(C)C</final\_answer>}\\[2pt]
\nodebadge{J}~N36 \emph{Despite the 3H vs.\,6H discrepancy, the overall
evidence supports the isopropyl ester}\\
\nodebadge{C}~N37 \emph{Agent commits despite unresolved discrepancy}

\end{description}

\medskip
\begin{center}
\begin{tikzpicture}[
  ncirc/.style={circle, draw, inner sep=1pt, minimum size=18pt,
                font=\small\sffamily\bfseries},
  nlbl/.style={font=\footnotesize, below=2pt},
]
\node[ncirc, fill=nodeE!25] (E) {E};
\node[nlbl] at (E.south) {N25};
\node[ncirc, fill=nodeJ!25, right=55pt of E] (J) {J};
\node[nlbl] at (J.south) {N33};
\node[ncirc, fill=nodeH!25, right=55pt of J] (H) {H};
\node[nlbl] at (H.south) {N30};
\node[ncirc, draw=gray!40, text=gray!40, right=55pt of H] (Hrev) {H$'$};
\draw[->, thick] (E) -- node[above, font=\footnotesize] {informs} (J);
\draw[->, thick, red!60!black] (J) -- node[above, font=\footnotesize] {contradicts} (H);
\draw[->, gray!40, dashed, thick] (H) -- node[above, font=\footnotesize, text=gray] {updates\_to} (Hrev);
\node[red!70, font=\bfseries\small] at ($(H)!0.5!(Hrev)+(0,-0.3)$) {\texttimes};
\end{tikzpicture}
\end{center}

\smallskip\noindent
\textbf{Annotator quote:}
\emph{\enquote{H-NMR: 1.46\,ppm (d,\,6H) — This is for isopropyl CH$_3$ groups,
but experimental shows 1.43\,ppm (d,\,3H)\,\ldots\ The experimental data
has only 3H in the doublet, not 6H.  So this isn't isopropyl ester.}}
— The agent \emph{explicitly} notes that the simulated $^{1}$H~NMR
contradicts its own hypothesis~(N30).  Yet no revised hypothesis is generated:
the \textsf{contradicts} edge from~N33 to~N30 has \textbf{no accompanying
\textsf{updates\_to}~edge}.  Instead, the agent dismisses the contradiction
as a simulation artefact and submits the isopropyl ester
(\texttt{CC(C)OC(=O)c1ccccc1N(C)C}) as the final answer~(N38).
\end{tcolorbox}


\definecolor{colENU}{RGB}{198,219,239}   % Evidence non-uptake — light blue
\definecolor{colUC}{RGB}{253,205,172}    % Untested claim — light orange
\definecolor{colFB}{RGB}{203,213,232}    % Fixed belief trace — light lavender
\definecolor{colCR}{RGB}{252,187,161}    % Contradiction without repair — light red

\definecolor{nodeE}{RGB}{86,146,194}     % Evidence node badge
\definecolor{nodeH}{RGB}{196,124,58}     % Hypothesis node badge
\definecolor{nodeJ}{RGB}{100,140,100}    % Judgment node badge
\definecolor{nodeT}{RGB}{136,136,136}    % Test node badge
\definecolor{nodeC}{RGB}{170,68,68}      % Commitment node badge
\definecolor{nodeF}{RGB}{80,80,80}       % Final-answer node badge

\newcommand{\nodebadge}[1]{%
  \tikz[baseline=(n.base)]{%
    \node[circle, draw, fill=node#1!25, inner sep=0pt,
          minimum size=9pt, font=\sffamily\bfseries\tiny] (n) {#1};%
  }%
}

\subsection{Illustrative traces}
\label{app:illustrative-traces}

For each major breakdown category, we include one representative trace excerpt
with the supporting quote the annotator used to assign the pattern.  For every
message in the excerpt, the epistemic nodes extracted by the annotator are
listed, together with the edges that link them.  A
\nodebadge{H}~denotes a \emph{hypothesis},
\nodebadge{T}~a \emph{test},
\nodebadge{E}~an \emph{evidence} observation,
\nodebadge{J}~a \emph{judgment}, and
\nodebadge{C}~a \emph{commitment} (premature answer).
Missing edges that \emph{should} appear in a sound reasoning chain are marked
with \textcolor{red!70}{\textsf{\bfseries\small$\times$}}.
The full annotated traces, with node and edge graphs, filtering by pattern,
and alternative candidate annotations, are browsable at
\url{https://lamalab-org.github.io/corral}.

\begin{tcolorbox}[
  colback=colENU!8, colframe=colENU!60!black,
  title={\small\sffamily\bfseries Evidence non-uptake},
  fontupper=\small,
  left=4pt, right=4pt, top=4pt, bottom=4pt,
  boxrule=0.4pt, arc=2pt,
  subtitle style={colback=colENU!25},
]
\textbf{Model:} Claude Sonnet 4.5\quad
\textbf{Environment:} NMR structure elucidation (spectra)\quad
\textbf{Scope:} 1\quad
\textbf{Trial:} 20

\tcblower
\begin{description}[font=\sffamily\bfseries\small, leftmargin=12pt, style=sameline, itemsep=6pt]

\item[Msg\,11 \normalfont\textsf{(Observation)}]
\texttt{obtain\_isomers\_from\_molecular\_formula}
$\to$~20~candidate SMILES returned for C$_{18}$H$_{14}$O$_2$, including
\texttt{COC(=O)/C=C\textbackslash{}C1=CC2=\allowbreak{}CC3=CC=CC=C3C=C2C=C1} (the correct answer)
.\\[2pt]
\nodebadge{T}~N13 \emph{Obtain isomers for molecular formula C$_{18}$H$_{14}$O$_2$}
$\xrightarrow{\textsf{observes}}$
\nodebadge{E}~N14 \emph{20 isomers returned, incl.\ several with methyl ester groups}

\item[Msg\,12 \normalfont\textsf{(Agent)}]
\enquote{\emph{I got 20 isomers with the molecular formula C$_{18}$H$_{14}$O$_2$.
Looking at the list, I can see several candidates.  Let me analyze which one
matches our NMR data\,\ldots}}\\
\textcolor{gray}{\footnotesize $\rightarrow$~calls \texttt{retrieve\_dbe\_formula}}\\[2pt]
No node extracted—the agent begins a tangential calculation instead of
using the isomer list.

\item[Msg\,14 \normalfont\textsf{(Agent)}]
\enquote{\emph{Let me recalculate the DBE for C$_{18}$H$_{14}$O$_2$\,\ldots\ Looking
back at the isomers list, let me focus on those with methyl ester
groups\,\ldots}}\\
\textcolor{gray}{\footnotesize $\rightarrow$~calls \texttt{simulate\_spectra} for methyl 9-phenanthrenecarboxylate}\\[2pt]
\nodebadge{T}~N15 \emph{Simulate spectra for methyl 9-phenanthrenecarboxylate
to compare with experimental}
\end{description}

\medskip
\begin{center}
\begin{tikzpicture}[
  ncirc/.style={circle, draw, inner sep=1pt, minimum size=18pt,
                font=\small\sffamily\bfseries},
  nlbl/.style={font=\footnotesize, below=2pt},
]
\node[ncirc, fill=nodeT!25] (T) {T};
\node[nlbl] at (T.south) {N13};
\node[ncirc, fill=nodeE!25, right=55pt of T] (E) {E};
\node[nlbl] at (E.south) {N14};
\node[ncirc, draw=gray!40, text=gray!40, right=55pt of E] (JH) {J/H};
\draw[->, thick] (T) -- node[above, font=\footnotesize] {observes} (E);
\draw[->, gray!40, dashed, thick] (E) -- node[above, font=\footnotesize, text=gray] {informs} (JH);
\node[red!70, font=\bfseries\small] at ($(E)!0.5!(JH)+(0,-0.3)$) {\texttimes};
\end{tikzpicture}
\end{center}

\smallskip\noindent
\textbf{Annotator quote:}
\emph{\enquote{\texttt{COC(=O)/C=C\textbackslash{}C1=CC2=\allowbreak{}CC3=CC=CC=C3C=C2C=C1},
\texttt{COC(=O)C12C=CC(\allowbreak{}C3=CC=CC=C31)C4=CC=CC=C24}}} — the 20 candidate
structures include the correct answer, yet node~N14 has \textbf{zero outgoing
\textsf{informs} edges}: the isomer list is never used to constrain any
judgment or hypothesis.  The agent instead guesses structures one by one
(phenanthrenecarboxylate, anthracenecarboxylate,~\ldots) without consulting
the evidence it already obtained.

\end{tcolorbox}

\begin{tcolorbox}[
  colback=colUC!8, colframe=colUC!60!black,
  title={\small\sffamily\bfseries Untested claim},
  fontupper=\small,
  left=4pt, right=4pt, top=4pt, bottom=4pt,
  boxrule=0.4pt, arc=2pt,
  subtitle style={colback=colUC!25},
]
\textbf{Model:} Claude Sonnet 4.5\quad
\textbf{Environment:} NMR structure elucidation (spectra)\quad
\textbf{Scope:} 1\quad
\textbf{Trial:} 22

\tcblower
\begin{description}[font=\sffamily\bfseries\small, leftmargin=12pt, style=sameline, itemsep=6pt]

\item[Msg\,3 \normalfont\textsf{(Observation)}]
\texttt{proton\_nmr\_spectra}
$\to$~\emph{\enquote{Deltas 2.23 (s, 3H), 2.81 (t, $J$\,=\,7.1\,Hz, 2H),
2.98 (t, $J$\,=\,7.1\,Hz, 2H), 7.17 (ddd, 7H), 7.90 (ddd, 2H).}}\\[2pt]
\nodebadge{T}~N1 \emph{Get proton NMR spectra to gather structural information}
$\xrightarrow{\textsf{observes}}$
\nodebadge{E}~N2 \emph{Proton NMR: 2.23\,(s,\,3H), 2.81\,(t,\,2H),
2.98\,(t,\,2H), 7.17\,(ddd,\,7H), 7.90\,(ddd,\,2H)}

\item[Msg\,4 \normalfont\textsf{(Agent)}]
\enquote{\emph{The two triplets at 2.81 and 2.98\,ppm suggest two CH$_2$ groups
coupled to each other (likely -CH$_2$-CH$_2$-).  The aromatic signals suggest a
substituted benzene ring.  The singlet at 2.23\,ppm suggests a methyl group
attached to a carbonyl (acetyl group).}}\\
\textcolor{gray}{\footnotesize $\rightarrow$~calls
\texttt{mass\_spectrometry\_spectra}}\\[2pt]
\nodebadge{E}~N2 $\xrightarrow{\textsf{informs}}$
\nodebadge{H}~N3 \emph{Two coupled CH$_2$ groups; singlet at 2.23\,ppm
suggests acetyl group; substituted benzene ring}

\item[Msg\,5 \normalfont\textsf{(Observation)}]
\texttt{mass\_spectrometry\_spectra} $\to$ \emph{\enquote{m/z 224.12 (intensity
100)\,\ldots}}\\[2pt]
The agent moves directly to mass spectrometry—no test is designed to
evaluate the CH$_2$-CH$_2$ or acetyl-group hypothesis.

\end{description}

\medskip
\begin{center}
\begin{tikzpicture}[
  ncirc/.style={circle, draw, inner sep=1pt, minimum size=18pt,
                font=\small\sffamily\bfseries},
  nlbl/.style={font=\footnotesize, below=2pt},
]
\node[ncirc, fill=nodeE!25] (E) {E};
\node[nlbl] at (E.south) {N2};
\node[ncirc, fill=nodeH!25, right=55pt of E] (H) {H};
\node[nlbl] at (H.south) {N3};
\node[ncirc, draw=gray!40, text=gray!40, right=55pt of H] (T) {T};
\draw[->, thick] (E) -- node[above, font=\footnotesize] {informs} (H);
\draw[->, gray!40, dashed, thick] (H) -- node[above, font=\footnotesize, text=gray] {tests} (T);
\node[red!70, font=\bfseries\small] at ($(H)!0.5!(T)+(0,-0.3)$) {\texttimes};
\end{tikzpicture}
\end{center}

\smallskip\noindent
\textbf{Annotator quote:}
\emph{\enquote{The two triplets at 2.81 and 2.98\,ppm suggest two CH$_2$ groups
coupled to each other (likely -CH$_2$-CH$_2$-).  The aromatic signals
suggest a substituted benzene ring.  The singlet at 2.23\,ppm suggests a
methyl group attached to a carbonyl (acetyl group).}}
— Node~N3 carries a detailed structural interpretation but has \textbf{zero
outgoing \textsf{tests} edges}.  The agent asserts the presence of an acetyl
group and coupled CH$_2$ groups as fact, then immediately requests mass
spectrometry to determine molecular weight rather than designing an
experiment (e.g., simulating $^{13}$C NMR for a candidate) to verify
the claim.

\end{tcolorbox}

\begin{tcolorbox}[
  colback=colFB!8, colframe=colFB!60!black,
  title={\small\sffamily\bfseries Fixed belief trace},
  fontupper=\small,
  left=4pt, right=4pt, top=4pt, bottom=4pt,
  boxrule=0.4pt, arc=2pt,
  subtitle style={colback=colFB!25},
]
\textbf{Model:} GPT-4o\quad
\textbf{Environment:} LAMMPS molecular dynamics (md)\quad
\textbf{Scope:} 2\quad
\textbf{Trial:} 33

\tcblower
\begin{description}[font=\sffamily\bfseries\small, leftmargin=12pt, style=sameline, itemsep=6pt]

\item[Msg\,15 \normalfont\textsf{(Observation)}]
\texttt{list\_files} $\to$ \emph{\enquote{\{\enquote{files}: [\enquote{/potentials/BKS},
\enquote{/potentials/EAM}, \enquote{/potentials/SW}, \enquote{/potentials/TERSOFF}]\}}}\\[2pt]
\nodebadge{T}~N10 \emph{List files in /potentials/ to find correct SW
potential path}
$\xrightarrow{\textsf{observes}}$
\nodebadge{E}~N11 \emph{Potentials directory contains BKS, EAM, SW, TERSOFF
subdirectories}

\item[Msg\,16 \normalfont\textsf{(Agent)}]
\enquote{\emph{The \texttt{/potentials/} directory contains a file named
\texttt{SW}, which likely corresponds to the Stillinger-Weber potential.
I will update the LAMMPS input script to use this file\,\ldots}}\\
\textcolor{gray}{\footnotesize $\rightarrow$~calls \texttt{write\_file}
to set \texttt{pair\_coeff * * /potentials/SW Si}}\\[2pt]
\nodebadge{E}~N11 $\xrightarrow{\textsf{informs}}$
\nodebadge{H}~N12 \emph{The SW potential file is located at /potentials/SW}

\item[] {\color{gray} $\vdots$ \footnotesize\itshape three error messages and re-writes later}

\item[Msg\,27 \normalfont\textsf{(Observation)}]
\texttt{run\_lammps} $\to$ \emph{\enquote{ERROR: Lost atoms: original 8
current 0}}\\[2pt]
\nodebadge{T}~N33 \emph{Run LAMMPS simulation with updated script using
/potentials/SW/Si.sw}
$\xrightarrow{\textsf{observes}}$
\nodebadge{E}~N34 \emph{LAMMPS simulation failed: Lost atoms}

\item[Msg\,28 \normalfont\textsf{(Agent)}]
\enquote{\emph{This might be due to the rapid heating rate or insufficient
equilibration.  I will adjust the heating rate and ensure proper
equilibration\,\ldots}}\\[2pt]
\nodebadge{E}~N34 \;(no \textsf{informs} to any existing H)\\
\nodebadge{H}~N35 \emph{Lost atoms error may be due to rapid heating rate
or insufficient equilibration}

\end{description}

\medskip
\begin{center}
\begin{tikzpicture}[
  ncirc/.style={circle, draw, inner sep=1pt, minimum size=18pt,
                font=\small\sffamily\bfseries},
  nlbl/.style={font=\footnotesize, below=2pt},
]
\node[ncirc, fill=nodeH!25] (H1) {H};
\node[nlbl] at (H1.south) {N12};
\node[ncirc, fill=nodeH!25, right=60pt of H1] (H2) {H};
\node[nlbl] at (H2.south) {N19};
\node[right=16pt of H2, font=\small, text=gray] {$\cdots$};
\node[ncirc, fill=nodeH!25, right=75pt of H2] (H5) {H};
\node[nlbl] at (H5.south) {N46};
\node[ncirc, fill=nodeH!25, right=60pt of H5] (H6) {H};
\node[nlbl] at (H6.south) {N57};
\draw[->, gray!40, dashed, thick] (H1) -- node[above, font=\footnotesize, text=gray] {updates\_to} (H2);
\node[red!70, font=\bfseries\small] at ($(H1)!0.5!(H2)+(0,-0.3)$) {\texttimes};
\draw[->, gray!40, dashed, thick] (H5) -- node[above, font=\footnotesize, text=gray] {updates\_to} (H6);
\node[red!70, font=\bfseries\small] at ($(H5)!0.5!(H6)+(0,-0.3)$) {\texttimes};
\end{tikzpicture}
\end{center}

\smallskip\noindent
\textbf{Annotator quotes:}\\
\noindent N12 ($\mathsf{msg\,16}$): \emph{\enquote{The \texttt{/potentials/}
directory contains a file named \texttt{SW}, which likely corresponds to the
Stillinger-Weber potential.}}\\
\noindent N35 ($\mathsf{msg\,28}$): \emph{\enquote{This might be due to the rapid
heating rate or insufficient equilibration.}}

\smallskip\noindent
The trace contains \textbf{6 hypothesis nodes and zero \textsf{updates\_to}
edges}.  The first hypothesis~(N12) mistakes the directory
\texttt{/potentials/SW} for a file, producing an error that is never traced
back to the belief.  After two more failures the agent introduces a second
independent hypothesis~(N35)—that lost atoms stem from heating rate—without
revising the earlier structural error.  Both beliefs persist unmodified
through the remaining 15~messages.

\end{tcolorbox}

\begin{tcolorbox}[
  colback=colCR!8, colframe=colCR!60!black,
  title={\small\sffamily\bfseries Contradiction without repair},
  fontupper=\small,
  left=4pt, right=4pt, top=4pt, bottom=4pt,
  boxrule=0.4pt, arc=2pt,
  subtitle style={colback=colCR!25},
]
\textbf{Model:} Claude Sonnet 4.5\quad
\textbf{Environment:} NMR structure elucidation (spectra)\quad
\textbf{Scope:} 2\quad
\textbf{Trial:} 45

\tcblower
\begin{description}[font=\sffamily\bfseries\small, leftmargin=12pt, style=sameline, itemsep=6pt]

\item[Msg\,17 \normalfont\textsf{(Observation)}]
\texttt{simulate\_spectra} for \texttt{CC(C)OC(=O)c1ccc\-cc1N(C)C}
$\to$~$^{1}$H NMR: 1.46\,(d,\,\textbf{6H}),\;2.97\,(s,\,6H),\;
4.75\,(p,\,1H),\;aromatic protons.\\[2pt]
\nodebadge{T}~N31 \emph{Simulate spectra for
\texttt{CC(C)OC(=O)c1ccccc1N(C)C}}
$\xrightarrow{\textsf{observes}}$
\nodebadge{E}~N25 \emph{Simulated: 1.46\,ppm (d,\,\textbf{6H});
experimental: 1.43\,ppm (d,\,\textbf{3H})}

\item[Msg\,18 \normalfont\textsf{(Agent)}]
\enquote{\emph{H-NMR: 1.46\,ppm (d,\,6H) — This is for isopropyl CH$_3$ groups, but
experimental shows 1.43\,ppm (d,\,3H).  The experimental data has only 3H in
the doublet, not 6H.  \textbf{So this isn't isopropyl ester.}  The
experimental data clearly shows CH(CH$_3$), not CH(CH$_3$)$_2$.}}\\
\textcolor{gray}{\footnotesize $\rightarrow$~calls
\texttt{get\_formula\_from\_smiles}}\\[2pt]
\nodebadge{E}~N25 $\xrightarrow{\textsf{informs}}$
\nodebadge{J}~N33 \emph{Simulated isopropyl ester shows 6H doublet at
1.46\,ppm, but experimental shows only 3H doublet at 1.43\,ppm}\\[2pt]
\nodebadge{J}~N33 $\xrightarrow{\textsf{contradicts}}$
\nodebadge{H}~N30 \emph{The compound is 2-(dimethylamino)benzoic acid
isopropyl ester}

\item[Msg\,20 \normalfont\textsf{(Agent --- final answer)}]
\enquote{\emph{Simulations are approximations and can have errors\,\ldots\
The key evidence: (1)~quartet at 5.05\,ppm, (2)~singlet at 2.69\,ppm
(6H) = N(CH$_3$)$_2$, (3)~4 aromatic protons, (4)~carbonyl at
168.30\,ppm.}}
\par\texttt{<final\_answer>CC(C)OC(=O)c1ccc\-cc1N(C)C</final\_answer>}\\[2pt]
\nodebadge{J}~N36 \emph{Despite the 3H vs.\,6H discrepancy, the overall
evidence supports the isopropyl ester}\\
\nodebadge{C}~N37 \emph{Agent commits despite unresolved discrepancy}

\end{description}

\medskip
\begin{center}
\begin{tikzpicture}[
  ncirc/.style={circle, draw, inner sep=1pt, minimum size=18pt,
                font=\small\sffamily\bfseries},
  nlbl/.style={font=\footnotesize, below=2pt},
]
\node[ncirc, fill=nodeE!25] (E) {E};
\node[nlbl] at (E.south) {N25};
\node[ncirc, fill=nodeJ!25, right=55pt of E] (J) {J};
\node[nlbl] at (J.south) {N33};
\node[ncirc, fill=nodeH!25, right=55pt of J] (H) {H};
\node[nlbl] at (H.south) {N30};
\node[ncirc, draw=gray!40, text=gray!40, right=55pt of H] (Hrev) {H$'$};
\draw[->, thick] (E) -- node[above, font=\footnotesize] {informs} (J);
\draw[->, thick, red!60!black] (J) -- node[above, font=\footnotesize] {contradicts} (H);
\draw[->, gray!40, dashed, thick] (H) -- node[above, font=\footnotesize, text=gray] {updates\_to} (Hrev);
\node[red!70, font=\bfseries\small] at ($(H)!0.5!(Hrev)+(0,-0.3)$) {\texttimes};
\end{tikzpicture}
\end{center}

\smallskip\noindent
\textbf{Annotator quote:}
\emph{\enquote{H-NMR: 1.46\,ppm (d,\,6H) — This is for isopropyl CH$_3$ groups,
but experimental shows 1.43\,ppm (d,\,3H)\,\ldots\ The experimental data
has only 3H in the doublet, not 6H.  So this isn't isopropyl ester.}}
— The agent \emph{explicitly} notes that the simulated $^{1}$H~NMR
contradicts its own hypothesis~(N30).  Yet no revised hypothesis is generated:
the \textsf{contradicts} edge from~N33 to~N30 has \textbf{no accompanying
\textsf{updates\_to}~edge}.  Instead, the agent dismisses the contradiction
as a simulation artefact and submits the isopropyl ester
(\texttt{CC(C)OC(=O)c1ccccc1N(C)C}) as the final answer~(N38).
\end{tcolorbox}
